\section{Formalizzazione dei casi d'uso significativi}\label{FormalizzazioneUseCase}
In questa sezione sono riportati i casi d'uso significativi dell'applicazione
con la loro formalizzazione secondo lo standard Cockburn.
In particolare viene riportato il caso d'uso riguardante la segnalazione di una nuova issue.

\subsection{Fully dressed use case}

\begin{table}[H]    
\small  
\def\arraystretch{1.3}

\begin{tabularx}{\linewidth}{|l|X|X|X|}

    \hline \textbf{Use Case} \#\ref{SegnalazioneUseCase} & 
    \multicolumn{3}{l|}{Creazione Issue} \\ 
    \hline \textbf{Goal In Context} & 
    \multicolumn{3}{p{10cm}|}{L'utente vuole creare una nuova issue} \\
    \hline \textbf{Preconditions} &
    \multicolumn{3}{p{10cm}|}{Avere ruolo di amministratore o utente} \\
    \hline \textbf{Success End Conditions} &
    \multicolumn{3}{p{10cm}|}{L'issue viene salvata e aggiunta alla lista delle issue} \\
    \hline \textbf{Failed End Conditions} &
    \multicolumn{3}{p{10cm}|}{L'issue non viene salvata} \\
    \hline \textbf{Primary Actor} &
    \multicolumn{3}{p{10cm}|}{Amministratore/utente} \\
    \hline \textbf{Trigger} & 
    \multicolumn{3}{p{10cm}|}{L'utente si autentica e clicca il pulsante "Crea issue"} \\
    \hline \multirow{5}{*}{\textbf{Main Scenario}} & \textbf{Step} & \textbf{User Action} & \textbf{System} \\
    \cline{2-4} & 1 & Clicca il bottone "Crea issue" & \\
    \cline{2-4} & 2 & & Mostra "Creazione Issue"\\
    \cline{2-4} & 3 & Compila i campi & \\
    \cline{2-4} & 4 &  & Mostra "Successo"\\
    \hline \multirow{5}{*}{\textbf{Extension}} & \textbf{Step} & \textbf{User Action} & \textbf{System} \\
    \cline{2-4} & 3a<Non compila tutti i campi> & Compila i campi & \\
    \cline{2-4} & 4a & & Mostra "Errore" \\
    \cline{2-4} & 5a & Clicca "OK" & \\
    \cline{2-4} & 6a & & Mostra "Creazione Issue" \\
    \hline \textbf{Open Issues} & \multicolumn{3}{l|}{Nessuna} \\
    \hline

\end{tabularx}

\end{table}
